\chapter{Multi-Dealer Secret Sharing} \label{chap:mdss}

\section{Overview}
In a secret sharing scheme \cite{Shamir79}, a dealer divides a secret into shares such that the secret remains hidden from anyone who holds an ``unauthorized'' set of shares, but can be recovered from any ``authorized'' set of shares.
These properties are referred to as privacy and correctness, respectively.

In this work, we consider a setting where {\em multiple} dealers may simultaneously distribute shares of their respective secrets.
We put forth a notion of {\em multi-dealer secret sharing} (MDSS) that achieves {\em stronger} security and correctness properties in this setting.

We name our security and correctness properties \emph{unlinkability} and \emph{multi-dealer} (MD) \emph{correctness} respectively, and give intuitive descriptions of them here.
We formally define an MDSS scheme and these properties in the following section.

\subsection{Unlinkability}
We propose a new security property for secret sharing that is stronger than the standard notion of privacy.
Intuitively, this property requires that given a set of shares from multiple dealers, an adversary cannot associate (i.e., link) any unauthorized subset of shares to a common dealer.

We present two definitions: unlinkability and one-more unlinkability. The first intuitively matches the  guarantee we would like to achieve, while the second is simpler and easier to use.
We will show that the two definitions are equivalent, and therefore that one can use the second without loss of generality.


\subsection{MD-Correctness}
asdf

\section{Definition}\label{sec:mdss-definition}
We study multi-dealer secret sharing for threshold access structures. Such a scheme is parameterized by four variables:

\begin{itemize}
	\item $\mdssRecThr$, the number of shares required to recover a secret.
	\item $\mdssPrivThr$, the number of shares one can emit before privacy and unlinkability are broken.
	\item $\mdssMaxDealers$, the number of sufficient dealers the scheme is able to tolerate.
	\item  $\mdssMaxShares$, the maximum number of shares that can be input to the reconstruction algorithm.
\end{itemize}

We consider the \emph{ramp} setting of Blakley and Meadows~\cite{C:BlaMea84}, where we allow for a gap between $\mdssRecThr$ and $\mdssPrivThr$ larger than $1$.
We now formally define multi-dealer secret sharing.

\begin{definition}[Multi-Dealer Secret Sharing]
	A $(\mdssRecThr, \mdssPrivThr, \mdssMaxDealers, \mdssMaxShares)$-multi-dealer secre sharing scheme (MDSS) with security parameter $\secpar$ is parameterized by $(\mdssRecThr, \mdssPrivThr, \mdssMaxDealers, \mdssMaxShares)$, each of which is an integer polynomial in $\secpar$, and is defined over a secret space $\mdssSecretSpace$ and a share space $\mdssShareSpace$.
	\hme{Technically share and secret space should also be families over the security parameter}

	An MDSS scheme consists of the following algorithms, which are expected to run in time polynomial in $\secpar$.

	\begin{itemize}
		\item $\mdssShareAlg(\secparam, \mdssSecret, \mdssNumShares) \rightarrow \set{\mdssShare{\mdssSecret}_i}_{i \in [\mdssNumShares]}$, takes as input a secret $\mdssSecret \in \mdssSecretSpace$ and an integer $\mdssNumShares$ and outputs a set of shares $\set{\mdssShare{\mdssSecret}_i}$ where each $\mdssShare{\mdssSecret}_i \in \mdssShareSpace$.
		\item $\mdssReconstructAlg(\mdssShare{\mdssSecret}_1, \dots, \mdssShare{\mdssSecret}_{\mdssWindowSize}) =: \set{\mdssShare{\mdssSecret}_i}$, takes as input shares $\mdssShare{\mdssSecret}_1$, \dots, $\mdssShare{\mdssSecret}_{\mdssWindowSize}$ where $\mdssWindowSize \leq \mdssMaxShares$ and outputs a (potentially empty) set of secrets.
	\end{itemize}
\end{definition}

We now define some notational shorthand.
For a secret $\mdssSecret$, an integer $\mdssNumShares$, and a list of unique indices $\indexSet \subseteq [\mdssNumShares]$, define the algorithm $\mdssProj(\mdssSecret, \mdssNumShares, \indexSet)$ as follows:

\procedureblock{$\mdssProj(\mdssSecret, \mdssNumShares, \indexSet)$}{
	\set{\mdssShare{\mdssSecret}_1, \dots, \mdssShare{\mdssSecret}_{\mdssNumShares}} \gets \mdssShareAlg(\mdssSecret, \mdssNumShares) \\
	\pcreturn \set{\mdssShare{\mdssSecret}_i}_{i \in \indexSet}
}

\begin{definition}[Unlinkability]
	\label{def:unlinkability}
	We say that an MDSS scheme is $\mdssULVar$\emph{-unlinkable} if for all admissible NUPPT adversaries $\adv = (\adv_1, \adv_2)$ and all $\secpar \in \NN$:
	\begin{equation*}
		\probsublong{
			\begin{gathered}
				\set{(\mdssSecret_i, (\mdssNumShares_{i, 0}, \indexSet_{i, 0}), (\mdssNumShares_{i, 1}, \indexSet_{i, 1}))}_{i \in [\mdssNumShareSets]}, \st \gets \adv_1(\secparam) \\
				b \sample \bin \\
				\forall i \in [\mdssNumShareSets]:\: \mdssShareSet_i \gets \mdssProj(\mdssSecret_i, \mdssNumShares_{i, b}, \indexSet_{i, b})
			\end{gathered}
		}{
			\adv_2(\st, \mdssShareSet_1, \dots, \mdssShareSet_{\mdssNumShareSets}) = b
		} \leq \frac{1}{2} + \mdssULVar(\secpar)
	\end{equation*}
	where $\adv$ is \emph{admissible} if it satitfies:
	\begin{enumerate}
		\item $\forall i \in [\mdssNumShareSets], \forall b \in \bin, |\indexSet_{i, b}| \leq \mdssPrivThr$
		\item $\indexSet_{0, 0}||\dots||\indexSet_{\mdssNumShareSets, 0} = \indexSet_{0, 1}||\dots||\indexSet_{\mdssNumShareSets, 1}$
	\end{enumerate}
\end{definition}
We refer to the above probability as the \emph{advantage} of $\adv$, and denote it by $\advantage{U}{\adv, \mathsf{MDSS}}$.
\hme{TODO: have a variable that denotes an MDSS scheme itself.}

While the above unlinkability definition is powerful, it is also somewhat cumbersome.
We therefore define a simpler definition for which it will be easier to prove that a scheme satisfies, and then show that the two definitions are in fact equivalent.

\begin{definition}[One-More Unlinkability]
	We say that an MDSS scheme is $\mdssULVar$\emph{-one-more-unlinkable} if for all admissible NUPPT adversaries $\adv = (\adv_1, \adv_2)$ and for all $\secpar \in \NN$:
	\begin{equation*}
		\begin{gathered}
			\probsublong{
				\begin{gathered}
					(\mdssSecret_0, \mdssSecret_1, \mdssNumShares_0, \mdssNumShares_1, \indexSet, i), \st \gets \adv_1(\secparam) \\
					\set{\mdssShare{\mdssSecret_0}_1, \dots, \mdssShare{\mdssSecret_0}_{\mdssNumShares_0 - 1}, \mdssShare{\mdssSecret_0}^*} \gets \mdssProj(\mdssSecret_0, \mdssNumShares_0, \indexSet)\\
					\set{\mdssShare{\mdssSecret_1}^*} \gets \mdssProj(\mdssSecret_1, \mdssNumShares_1, i) \\
				\end{gathered}
			}{
				\adv_2(\st, \mdssShare{\mdssSecret_0}_1, \dots, \mdssShare{\mdssSecret_0}_{\mdssNumShares_0 - 1}, \mdssShare{\mdssSecret_b}^*) = b
			} \\
			\leq \frac{1}{2} + \mdssULVar(\secpar)
		\end{gathered}
	\end{equation*}
\end{definition}
where $\adv$ is \emph{admissible} if it is the case that $|\indexSet| \leq \mdssPrivThr$.
We refer to the above probability as the \emph{advantage} of $\adv$, and denote it by $\advantage{OMU}{\adv, \mathsf{MDSS}}$.

\begin{definition}[MD-Correctness]
	We say that a $(\mdssRecThr, \mdssPrivThr, \mdssMaxDealers, \mdssMaxShares)$-MDSS scheme is $\mdssCVar$\emph{-correct} if for all $\secpar \in \NN$ and for all $\set{(\mdssSecret_i, \mdssNumShares_i, \indexSet_i)}_{[\mdssNumShareSets]}$ where the following hold:
	\begin{itemize}
		\item $\forall i \in [\mdssNumShareSets],\: \indexSet_i \subseteq [\mdssNumShares_i]$

		\item $\sum_{i \in [\mdssNumShareSets]} |\indexSet_i| \leq \mdssMaxShares$
		\item $|\set{\indexSet_i \text{ s.t. } |\indexSet_i| \geq \mdssRecThr}| \leq \mdssMaxDealers$

		\item $\bigcap_{i, j \in [\mdssNumShareSets],\: i \neq j} \indexSet_i \cap \indexSet_j = \emptyset$
	\end{itemize}
	it holds that:
	\begin{equation*}
		\probsublong{
			\mdssShareSet := \bigcup_{i \in [\mdssNumShareSets]} \mdssProj(\mdssSecret_i, \mdssNumShares_i, \indexSet_i)
		}{
			\mdssReconstructAlg(\mdssShareSet) \neq \set{\mdssSecret_i \text{ s.t. } |\indexSet_i| \geq \mdssRecThr}
		} \leq \mdssCVar(\secpar)
	\end{equation*}
\end{definition}


\section{Proofs}

We prove here that unlinkability is a strictly stronger requirement than traditional secret-sharing privacy, as well as that unlinkability and one-more unlinkability are equivalent.

\subsection{Unlinkability vs. Privacy}
We begin by recalling the standard notion of privacy for secret sharing, which states that any unauthorized set of shares hides the secret.

\begin{definition}[Privacy]
	We way that a $(\mdssRecThr, \mdssPrivThr, \mdssMaxDealers, \mdssMaxShares)$-MDSS scheme with secret space $\mdssSecretSpace$ is $\mdssULVar$\emph{-private} if for any $s_0, s_1 \in \mdssSecretSpace$ and any $\mdssNumShares \in \ZZ$ where $\mdssNumShares \leq \mdssPrivThr$ it holds that:
	\begin{equation*}
		\probsublong{
			\begin{gathered}
				\mdssShareSet_0 \gets \mdssShareAlg(\secparam, \mdssNumShares, \mdssSecret_0) \\
				\mdssShareSet_1 \gets \mdssShareAlg(\secparam, \mdssNumShares, \mdssSecret_1)
			\end{gathered}
		} {
			\adv(S_b) = b
		} \leq \frac{1}{2} + \mdssULVar(\secpar)
	\end{equation*}
\end{definition}

We now show that unlinkability is strictly stronger than privacy.

\begin{lemma}
	If a $(\mdssRecThr, \mdssPrivThr, \mdssMaxDealers, \mdssMaxShares)$-MDSS scheme is $\mdssULVar$-unlinkable then it is $\mdssULVar$-private.
\end{lemma}
\begin{proof}
	Consider an unlinkability adversary that outputs the following in the first step of the game:
	$$
		\set{(\mdssSecret_0, (\mdssNumShares, [\mdssNumShares]), (0, \emptyset)), (\mdssSecret_1, (0, \emptyset), (\mdssNumShares, [\mdssNumShares]))}
	$$
	Such an adversary must distinguish between $\mdssNumShares$ shares of $\mdssSecret_0$ and $\mdssNumShares$ of $\mdssSecret_1$, which corresponds exactly to the privacy game.
	\hme{Change to full contrapositive proof}
\end{proof}

\begin{lemma}
	Define the function $\mdssULVar^*(\secpar) = \frac{1}{2} - 2^{-\secpar}$.
	% There exists a $\mdssULVar^0$-private secret sharing scheme that is not $\mdssULVar^*$-unlinkable.
	For any $\mdssULVar$-private secret sharing scheme there exists a scheme that is $\mdssULVar$-private and \emph{not} $\mdssULVar^*$-unlinkable.
\end{lemma}

\begin{proof}
	Consider a secret-sharing scheme with sharing algorithm $\mdssShareAlg$ and share space $\mdssShareSpace$.
	Next consider a secret-sharing scheme with the modified sharing algorithm $\mdssShareAlg'$ and share space $\mdssShareSpace \times \bin^{\secpar}$.
	\procedureblock{$\mdssShareAlg'(\secparam, \mdssNumShares, \mdssSecret)$}{
		\set{\mdssShare{\mdssSecret}_0, \dots, \mdssShare{\mdssSecret}_{\mdssNumShares}} \gets \mdssShareAlg(\secparam, \mdssNumShares, \mdssSecret) \\
		r \sample \bin^{\secpar} \\
		\pcreturn \set{(\mdssShare{\mdssSecret}_0, r), \dots, (\mdssShare{\mdssSecret}_{\mdssNumShares}, r)}
	}

	Clearly the scheme remains $\mdssULVar^*$-private as the $r$ values contain no information about $s$.

	Next, consider the following one-more unlinkability adversary $\adv = (\adv_1, \adv_2)$
	\begin{pchstack}[center, space=1em]
		\procedure{$\adv_1(\secparam)$}{
			\pcreturn (s_0, s_1, 2, 1, \set{1}, \set{1}), \bot
		}

		\procedure{$\adv_2(\st, (\mdssBlindShare_0, r_0), (\mdssBlindShare_1, r_1))$}{
			\pcif r_0 = r_1 \pcthen \pcreturn 0 \\
			\pcelse \pcreturn 1
		}
	\end{pchstack}
	Clearly, $\adv$ only loses the one-more unlinkability game when $b = 1$ and $r_0 = r_1$.
	This happens with probability $2^{-(\secpar + 1)}$, and so the scheme is not $\frac{1}{2} - 2^{-\secpar}$ unlinkable.
\end{proof}

\subsection{Unlinkability vs. One-More Unlinkability}

We show here that unlinkability and one-more unlinkability are equivalent definitions.
First, it should be clear to see that any scheme satisfying unlinkability also satisfies one-more unlinkability.
Therefore, we only prove the relationship in the other direction.

\begin{theorem}
	If a MDSS scheme satisfies $\negl[]$-one-more unlinkability,  then it satisfies $\negl[]$-unlinkability.
\end{theorem}

\begin{proof}
	For game $\game_i$, let $\G_i$ denote the probability that the adversary successfully guesses the bit in the game.

	Recall that the adversary $\adv$ queries on $\set{(\mdssSecret_i, (\mdssNumShares_{i, 0},\indexSet_{i, 0}), (\mdssNumShares_{i, 1}, \indexSet_{i, 1}))}_{i \in [\mdssNumShareSets]}$.
	Let $q(\secpar) = \sum_i |\indexSet_{i, 0}| = \sum_i |\indexSet_{i, 1}|$.
	We now consider a sequence of games.
	Let $\game_0$ denote the original unlinkability game, and $\game_{i}$ denote the game where shares at indices $\set{q(\secpar) - i + 1, q(\secpar) - i + 2, \dots, q(\secpar)}$ are replaced with $\mdssShareAlg(s, 1)$, where $s$ is sampled uniformly from $\mdssShareSpace$.
	Note that $\game_0$ is identical to the unlinkability game described in \cref{def:unlinkability}.

	\begin{claim}
		For all $j \in \set{0, 1, \dots, q(\secpar) - 1}$, for all admissible adversaries $\adv$, $\exists$ a negligible function $\negl[]$ such that $|\G_j -\G_{j+1}| \leq \negl$.
	\end{claim}

	\begin{proof}
		Recall that in its first phase the unlinkability adversary $\adv$ outputs $\set{(\mdssSecret_i, (\mdssNumShares_{i, 0},\indexSet_{i, 0}), (\mdssNumShares_{i, 1}, \indexSet_{i, 1}))}_{i \in [\mdssNumShareSets]}$.
		For $\game_j$ we define the following notation.
		\begin{itemize}
			\item Let $\mdssRelBatch$ be the smallest integer such that $\sum_{i \in [\mdssRelBatch]} |\indexSet_i| \geq j$. In other words, $\mdssProj(\mdssSecret, \mdssNumShares_{\mdssRelBatch, b}, \indexSet_{\mdssRelBatch, b})$ is the ``batch'' in which the share $\mdssBlindShare_j$ appears.
			\item Let $\mdssLIBRB = \sum_{i \in [\mdssRelBatch - 1]} |\indexSet_i|$, i.e.\ the last index before the relevant batch.
			\item Let $\mdssIWRB = j - \mdssLIBRB$, i.e.\ the index within the relevant batch at which $\mdssBlindShare_j$ appears.
			\item Define the function $\mdssSampleShare(\secparam)$ as one that samples $\mdssSecret \sample \mdssSecretSpace$ and returns $\mdssShareAlg(\secparam, \mdssSecret, 1)$.
			      Similarly, define $\mdssSampleShare^w(\secparam)$ to be the function that computes $\mdssBlindShare_i = \mdssSampleShare(\secparam)$ for $i \in [w]$ and outputs $\set{\mdssBlindShare_i}_{i \in [w]}$.
		\end{itemize}

		Now consider the following one-more unlinkability adversary $\bdv^j = (\bdv^j_1, \bdv^j_2)$.
		\begin{pcvstack}[center, space=1em]
			\procedure{$\bdv^j_1(\secparam)$}{
				\set{(\mdssSecret_i, (\mdssNumShares_{i, 0},\indexSet_{i, 0}), (\mdssNumShares_{i, 1}, \indexSet_{i, 1}))}_{i \in [\mdssNumShareSets]}, \st \gets \adv_1(\secparam) \\
				b \sample \bin \\
				\forall i \in [\mdssRelBatch - 1]: \mdssShareSet_i \gets \mdssProj(\mdssSecret_i, \mdssNumShares_{i, b}, \indexSet_{i, b}) \\
				\forall i \in [\mdssRelBatch + 1, \mdssNumShareSets]: \mdssShareSet_i \gets \mdssSampleShare^{\mdssNumShares_{i, b}}(\secparam) \\
				\mdssSecret' \sample \mdssSecretSpace \\
				\pcreturn (\mdssSecret_j, \mdssSecret', \mdssIWRB, 1, 1), (\st, b, \mdssNumShares_{k, b}, \set{\mdssShareSet_i}_{i \in [\mdssRelBatch]}, \set{\mdssShareSet_i}_{i \in [j+1, \mdssNumShareSets]})
			}

			\procedure{$\bdv_2^j((\st, b, \mdssNumShares_{\mdssRelBatch, b}, \set{\mdssShareSet_i}_{i \in [\mdssRelBatch]}, \set{\mdssShareSet_i}_{i \in [j+1, \mdssNumShareSets]}), \mdssBlindShare_1, \dots, \mdssBlindShare_{\mdssIWRB})$}{
				\forall k \in [\mdssIWRB + 1, \mdssNumShares_{\mdssRelBatch, b}]: \mdssBlindShare_k \gets \mdssSampleShare(\secparam) \\
				\mdssShareSet_j := \set{\mdssBlindShare_k}_{k \in [\mdssNumShares_{\mdssRelBatch, b}]} \\
				b' \gets \adv_2(\st, \set{\mdssShareSet_i}_{i \in [\mdssNumShareSets]}) \\
				\pcif b' = b \pcthen \pcreturn 0 \\
				\pcelse \pcreturn 1
			}
		\end{pcvstack}
		$\bdv^j$ invokes $\adv$, flips its own bit $b$, and then simulates all shares that do not correspond to the relevant batch $\mdssRelBatch$.
		$\bdv^j$ then receives $\mdssIWRB$ shares of the relevant batch (the last one of which may have been replaced with a random share), and simulates the remaining shares.
		$\bdv^j$ then invokes $\adv^j$ on the aggregated share sets, and checks whether $\adv$ succeeds in guessing $b$.

		We can now analyze $\bdv^j$'s advantage.
		\begin{align*}
			\advantage{omu}{\bdv, \mathsf{MDSS}} & = \prob{b' = b} \\
			                                     & =
		\end{align*}



		% $\bdv^j$ begins by invoking $\adv_1(\secparam)$ and receiving $\set{(\mdssSecret_i, (\mdssNumShares_{i, 0},\indexSet_{i, 0}), (\mdssNumShares_{i, 1}, \indexSet_{i, 1}))}_{i \in [\mdssNumShareSets]}$.
		% $\bdv$ samples $b \sample \bin$ and simulates all shares $\mdssBlindShare_1, \dots, \mdssBlindShare_{\ell}$ accordingly.
		% $\bdv$ then samples $\mdssSecret'$ and outputs $(\mdssSecret_j, \mdssSecret', m, 1, 1)$, receiving $(\mdssBlindShare_{\ell + 1}, \dots, \mdssBlindShare_j)$.
		% Finally, $\bdv$ simulates $\mdssBlindShare_{j+1},\dots,\mdssBlindShare_{q(\secpar)}$ each by computing a share of a sampled secret.
	\end{proof}


\end{proof}
